\subsection{TENSOR PRODUCT OF A FINITE FAMILY OF ALGEBRAS}

A always denotes a commutative ring with unit element. Let \((\mathbf{E}_i)_{i\in I}\) be a finite family of A-algebras and let \(\mathbf{E} = \bigotimes_{i\in I}\mathbf{E}_i\) be the tensor product A-module of the A-modules \(\mathbf{E}_i\) (II, § 3, no. 9). We shall define an A-algebra structure on E. Let \(m_i\colon \mathbf{E}_i\otimes_\mathbf{A}\mathbf{E}_i\to \mathbf{E}_i\) be the A-linear mapping defining the multiplication on \(\mathbf{E}_i\) (§ 1, no. 3). Consider the A-linear mapping

\[
m ^ {\prime} = \bigotimes_ {i \in I} m _ {i}: \bigotimes_ {i \in I} (E _ {i} \otimes_ {A} E _ {i}) \rightarrow \bigotimes_ {i \in I} E _ {i} = E;
\]

the composite mapping

\[
\left(\bigotimes_ {i \in I} \mathrm{E} _ {i}\right) \otimes_ {\mathrm{A}} \left(\bigotimes_ {i \in I} \mathrm{E} _ {i}\right) \xrightarrow {\tau} \bigotimes_ {i \in I} \left(\mathrm{E} _ {i} \otimes_ {\mathrm{A}} \mathrm{E} _ {i}\right) \xrightarrow {m ^ {\prime}} \bigotimes_ {i \in I} \mathrm{E} _ {i}
\]

where \(\tau\) is the associativity isomorphism (II, §3, no. 9) is an A-linear mapping \(m\colon \mathbf{E} \otimes_{\mathbf{A}} \mathbf{E} \to \mathbf{E}\) ; we shall see that \(m\) defines an (associative and unital) algebra structure on \(\mathbf{E}\) . For, on explicitly performing the multiplication defined by \(m\) , we obtain the formula

\[
\left(\bigotimes_ {i \in I} x _ {i}\right) \left(\bigotimes_ {i \in I} y _ {i}\right) = \bigotimes_ {i \in I} \left(x _ {i} y _ {i}\right) \quad \text { for } x _ {i}, y _ {i} \text { in } E _ {i} \text { and } i \in I.\tag{1}
\]

It is therefore seen already, by linearity, that if \(e_i\) is the unit element of \(\mathbf{E}_i\) , \(e = \bigotimes_{i \in I} e_i\) is unit element of \(\mathbf{E}\) . On the other hand, the associativity of each of the \(\mathbf{E}_i\) implies the relation

\[
\left(\left(\bigotimes_ {i \in I} x _ {i}\right) \left(\bigotimes_ {i \in I} y _ {i}\right)\right) \left(\bigotimes_ {i \in I} z _ {i}\right) = \bigotimes_ {i \in I} \left(x _ {i} y _ {i} z _ {i}\right) = \left(\bigotimes_ {i \in I} x _ {i}\right) \left(\left(\bigotimes_ {i \in I} y _ {i}\right) \left(\bigotimes_ {i \in I} z _ {i}\right)\right)
\]

whence, by linearity, the relation \(x(yz) = (xy)z\) for all x, y, z in E.

DEFINITION 1. Given a family \((\mathbf{E}_i)_{i\in I}\) of algebras over A, the tensor product of this family, denoted by \(\bigotimes_{i\in I}\mathbf{E}_i\) (or, when I is the interval \([1,n]\) of N, \(\mathbf{E}_1\otimes_\mathbf{A}\mathbf{E}_2\otimes \dots \otimes_\mathbf{A}\mathbf{E}_n\) , or simply \(\mathbf{E}_1\otimes \mathbf{E}_2\otimes \dots \otimes \mathbf{E}_n\) ) is the algebra obtained by giving the tensor product of the A-modules \(\mathbf{E}_i\) the multiplication defined by (1).

Relation (1) shows that the tensor product \(\bigotimes_{i\in I}E_i^0\) of the opposite algebras to the \(E_{i}\) is the opposite algebra to \(\bigotimes_{i\in I}E_{i}\) ; in particular, if the \(E_{i}\) are commutative, so is \(\bigotimes_{i\in I}E_{i}\) .

Let \((\mathbf{E}_i)_{i\in I}\) and \((\mathbf{F}_i)_{i\in I}\) be two families of A-algebras with the same finite indexing set I. For each \(i\in I\) , let \(f_{i}\colon \mathbf{E}_{i}\to \mathbf{F}_{i}\) be an A-algebra homomorphism. Then the A-linear mapping

\[
f = \bigotimes_ {i \in I} f _ {i}: \bigotimes_ {i \in I} E _ {i} \rightarrow \bigotimes_ {i \in I} F _ {i}
\]

is an A-algebra homomorphism, as follows from (1).

For every partition \((\mathbf{I}_{j})_{j\in J}\) of I, the associativity isomorphisms

\[
\bigotimes_ {j \in J} \left(\bigotimes_ {i \in I _ {j}} \mathrm{E} _ {i}\right)\rightarrow \bigotimes_ {i \in I} \mathrm{E} _ {i}
\]

(II, § 3, no. 9) are also algebra isomorphisms, as follows from (1) and their definitions.

When I is the interval \([1, n]\) of \(\mathbf{N}\) and all the algebras \(\mathbf{E}_i\) are equal to the same algebra \(\mathbf{E}\) , the tensor product algebra \(\bigotimes_{i \in I} \mathbf{E}_i\) is also denoted by \(\mathbf{E}^{\otimes n}\) .

We shall restrict our attention in the remainder of this no. to the properties of tensor products of two algebras, leaving to the reader the task of extending them to tensor products of arbitrary finite families.

Let E, F be two A-algebras; if a (resp. b) is a left ideal of E (resp. F), the canonical image \(\operatorname{Im}(a \otimes b)\) of \(a \otimes_{A} b\) in \(E \otimes_{A} F\) is a left ideal of \(E \otimes_{A} F\) ; there are analogous results when ``left ideal'' is replaced by ``right ideal'' or ``two-sided ideal''. Moreover:

PROPOSITION 1. Let E, F be two A-algebras and a (resp. b) a two-sided ideal of E (resp. F). Then the canonical A-module isomorphism

\[
(\mathrm{E} / \mathfrak {a}) \otimes (\mathrm{E} / \mathfrak {b}) \rightarrow (\mathrm{E} \otimes \mathrm{F}) / (\operatorname{Im} (\mathfrak {a} \otimes \mathrm{F}) + \operatorname{Im} (\mathrm{E} \otimes \mathfrak {b}))
\]

(II, § 3, no. 6, Corollary 1 to Proposition 6) is an algebra isomorphism.

This follows from (1) and the definition given loc. cit.

COROLLARY 1. Let E be an A-algebra and a an ideal of A. Then the A-module aE is a two-sided ideal of E and the canonical (A/a)-module isomorphism

\[
(\mathbf {A} / \mathfrak {a}) \otimes_ {\mathbf {A}} \mathbf {E} \rightarrow \mathbf {E} / \mathfrak {a} \mathbf {E}
\]

is an (A/a)-algebra isomorphism.

COROLLARY 2. If \(\mathfrak{a},\mathfrak{b}\) are two ideals of \(\mathbf{A}\) , \((\mathbf{A} / \mathfrak{a})\otimes_{\mathbf{A}}(\mathbf{A} / \mathfrak{b})\) is canonically isomorphic to \(\mathbf{A} / (\mathfrak{a} + \mathfrak{b})\) .

COROLLARY 3. Let E, F be two A-algebras and a an ideal of A contained in the annihilator of F. Then the (A/a)-algebra \(E \otimes_{A} F\) is canonically isomorphic to (E/aE) \(\otimes_{A/a} F\) .

PROPOSITION 2. Let \((\mathbf{E}_{\lambda})_{\lambda \in L}\) and \((\mathbf{F}_{\mu})_{\mu \in M}\) be two families of A-algebras. The canonical mapping (II, §3, no. 7)

\[
\left(\bigoplus_ {\lambda \in \mathbf {L}} \mathrm{E} _ {\lambda}\right) \otimes_ {\mathrm{A}} \left(\bigoplus_ {\mu \in \mathrm{M}} \mathrm{F} _ {\mu}\right)\rightarrow \bigoplus_ {(\lambda , \mu) \in \mathrm{L} \times \mathrm{M}} \left(\mathrm{E} _ {\lambda} \otimes_ {\mathrm{A}} \mathrm{F} _ {\mu}\right)
\]

is an algebra isomorphism.

This follows immediately from II, § 3, no. 7, Proposition 7 and the definition of multiplication on \(\mathbf{E} \otimes \mathbf{F}\) .

PROPOSITION 3. Let A, B be two commutative rings, \(\rho: \mathbf{A} \to \mathbf{B}\) a ring homomorphism and E, F two A-algebras. Then the canonical B-module isomorphism

\[
\rho^ {*} (\mathrm{E}) \otimes_ {\mathbf {B}} \rho^ {*} (\mathrm{F}) \rightarrow \rho^ {*} (\mathrm{E} \otimes_ {\mathbf {A}} \mathrm{F})
\]

(II, § 5, no. 1, Proposition 3) is a B-algebra isomorphism.

PROPOSITION 4. Let A, B be two commutative rings, \(\rho: \mathbf{A} \to \mathbf{B}\) a ring homomorphism, E an A-algebra and F a B-algebra. Then the canonical A-module isomorphism

\[
\rho_ {*} (\mathbf {F}) \otimes_ {\mathbf {A}} \mathbf {E} \rightarrow \rho_ {*} (\mathbf {F} \otimes_ {\mathbf {B}} \rho^ {*} (\mathbf {E}))
\]

(II, § 5, no. 2, Proposition 6) is an A-algebra isomorphism.

The verifications are trivial on account of § 1, no. 5.

In particular, the A-algebra structure on B \(\otimes_{\mathbf{A}} \mathbf{E}\) , obtained by restricting the ring B of scalars to A, is identical with the structure of the algebra B \(\otimes_{\mathbf{A}} \mathbf{E}\) , the tensor product of the A-algebras B and E.

Finally, if \((\mathbf{A}_i, \phi_{ji})\) is a direct system of commutative rings, \((\mathbf{E}_i, f_{ji})\) and \((\mathbf{F}_i, g_{ji})\) two direct systems of \(\mathbf{A}_i\) -algebras (§ 1, no. 6) and \(\mathbf{A} = \varinjlim \mathbf{A}_i\) , the canonical A-module isomorphism

\[
\varinjlim \left(\mathrm{E} _ {i} \otimes_ {\mathrm{A} _ {i}} \mathrm{F} _ {i}\right)\rightarrow (\varinjlim \mathrm{E} _ {i}) \otimes_ {\mathrm{A}} (\varinjlim \mathrm{F} _ {i})
\]

(II, § 6, no. 3, Proposition 7) is also an A-algebra isomorphism, as follows from the definitions.

Examples of tensor products of algebras. (1) Let A be a commutative ring and M, N two A-modules; the canonical mapping

\[
\operatorname{End} _ {\mathbf {A}} (\mathbf {M}) \otimes_ {\mathbf {A}} \operatorname{End} _ {\mathbf {A}} (\mathbf {N}) \rightarrow \operatorname{End} _ {\mathbf {A}} (\mathbf {M} \otimes_ {\mathbf {A}} \mathbf {N})\tag{2}
\]

(II, §4, no. 4) is an A-algebra homomorphism, as follows from II, §3, no. 2, formula (5). When M or N is a finitely generated projective A-module, we know that this homomorphism is bijective (II, §4, no. 4, Proposition 4). In particular we recover the definition of the tensor product of two square matrices.

\begin{enumerate}
\def\labelenumi{(\arabic{enumi})}
\setcounter{enumi}{1}
\tightlist
\item
  Let \(S, T\) be two monoids and \(A^{(S)}\) and \(A^{(T)}\) the algebras of the monoids \(S\) and \(T\) over the ring \(A\) (III, § 2, no. 6); then there is a canonical \(A\) -algebra isomorphism
\end{enumerate}

\[
\mathbf {A} ^ {(\mathbf {S})} \otimes_ {\mathbf {A}} \mathbf {A} ^ {(\mathbf {T})} \rightarrow \mathbf {A} ^ {(\mathbf {S} \times \mathbf {T})}.\tag{3}
\]

The elements \(e_s \otimes e_t\) (resp. \(e_{(s,t)}\) ), where \(s\) runs through \(S\) and \(t\) runs through \(T\) , form a basis of \(A^{(S)} \otimes_A A^{(T)}\) by virtue of II, § 3, no. 7, Corollary 2 to Proposition 7 (resp. of \(A^{(S \times T)}\) ); the desired isomorphism is obtained by mapping \(e_s \otimes e_t\) to \(e_{(s,t)}\) and it follows from the definitions that this is indeed an algebra isomorphism.

